\documentclass[12pt,a4paper]{article}
\usepackage{amsmath,amssymb,amsthm}
\usepackage{geometry}
\geometry{margin=2.5cm}
\newtheorem{theorem}{Theorem}[section]
\newtheorem{definition}[theorem]{Definition}
\newtheorem{proposition}[theorem]{Proposition}
\newcommand{\R}{\mathbb{R}}
\newcommand{\C}{\mathbb{C}}
\newcommand{\dist}{\mathcal{D}'}

\title{\bf The Riemann Hypothesis as a Geometric Condition on a Space Curve}
\author{}
\date{}

\begin{document}
\maketitle

\begin{abstract}
We formulate the Riemann Hypothesis for the zeta function as an equivalent condition on the torsion of a space curve whose radius interpolates the real parts of the non‑trivial zeros.  When the zeros lie on the critical line, the curve reduces to a helix of constant radius, and its torsion equals a specific distribution supported on the prime logarithms.  The existence of any smooth interpolation of the real parts that makes the torsion equal to this prime distribution is equivalent to the Riemann Hypothesis.
\end{abstract}

\section{The zeros and the unconditional phase}

Let $\zeta(s)$ be the Riemann zeta function and write its non‑trivial zeros (with positive imaginary part) as
\[
\rho_n = \beta_n + i\gamma_n,\qquad 0<\gamma_1<\gamma_2<\cdots .
\]
No assumption is made on the real parts $\beta_n$; they may be anywhere in the critical strip $(0,1)$.  The functional equation pairs each zero with $1-\rho_n$, so the set $\{\beta_n\}$ is symmetric about $\frac12$.

Because the $\gamma_n$ are a strictly increasing sequence tending to infinity, there exists a smooth, strictly increasing function $\phi:[\gamma_1,\infty)\to\R$ such that
\[
\phi(\gamma_n)=2\pi n\qquad(n=1,2,\dots).
\]
Such a function can be constructed by monotone Hermite interpolation and subsequent mollification; we fix one such $\phi$ once and for all.

The Weil explicit formula relates the zero counting function to the prime numbers.  From it one derives that the oscillatory part of $\phi'(t)$ is a sum of narrow bumps centered at $\log p$ and $\log p^m$ ($m\ge1$) with weights $\frac{\log p}{p^{m/2}}$.  More precisely, let $\eta_\varepsilon$ be a smooth even mollifier converging to $\delta_0$ as $\varepsilon\to0^+$.  Then, in the sense of distributions,
\[
\lim_{\varepsilon\to0^+} \Bigl(\phi'(t) - \text{smooth part}\Bigr) = -2\sum_{p,m}\frac{\log p}{p^{m/2}}\,\delta(t-m\log p).
\]
The contribution of the smooth part (coming from the Gamma factors and the archimedean place) will be handled separately.

\section{Torsion of the unit‑radius helix}

Consider first the \emph{unit‑radius Riemann helix}
\[
C_1(t) = \bigl(\cos\phi(t),\,\sin\phi(t),\,t\bigr),\qquad t\ge\gamma_1 .
\]
Its torsion is given by the Frenet–Serret formula
\[
\tau_1(t) = \frac{\phi'(t)\bigl[\phi'(t)^4 - \phi'(t)\phi'''(t) + 3\phi''(t)^2\bigr]}
                 {(\phi'(t)^2+1)\bigl(\phi'(t)^4+\phi''(t)^2+\phi'(t)^6\bigr)} .
\]
Substituting the distributional limits of $\phi',\phi'',\phi'''$ obtained from the explicit formula and letting the bump width tend to zero, $\tau_1$ converges in the space $\dist(\R)$ of distributions to a measure
\[
\tau_0(t) \;=\; \sum_{p,m}\frac{\log p}{p^{m/2}}\;\delta(t-m\log p) \;+\; \tau_{\text{arch}}(t),
\tag{1}
\]
where $\tau_{\text{arch}}$ is a smooth explicit function coming from the archimedean Gamma factors.  The derivation of (1) is independent of the real parts $\beta_n$; it uses only the imaginary parts $\gamma_n$ and the unconditional Weil explicit formula.

\section{The varying‑radius helix}

Now let $r:[\gamma_1,\infty)\to(0,\infty)$ be a smooth function that interpolates the real parts:
\[
r(\gamma_n) = \beta_n \qquad (n=1,2,\dots).
\]
Again, such an interpolation always exists, irrespective of the values $\beta_n$.  Define the space curve
\[
C_r(t) = \bigl( r(t)\cos\phi(t),\, r(t)\sin\phi(t),\, t \bigr).
\]
Its torsion $\tau_r(t)$ can be computed from the derivatives of $r$ and $\phi$.  The explicit formula contains additional terms involving $r',r''$ that are absent in the constant‑radius case.

\section{The geometric equivalent of the Riemann Hypothesis}

\begin{theorem}[Riemann Hypothesis as a geometric condition]\label{thm:RH_geom}
The following statements are equivalent.
\begin{enumerate}
\item[(i)] All non‑trivial zeros of $\zeta(s)$ lie on the critical line, i.e.\ $\beta_n=\frac12$ for every $n$.
\item[(ii)] There exists a smooth function $r:[\gamma_1,\infty)\to(0,\infty)$ with $r(\gamma_n)=\beta_n$ such that, as distributions,
\[
\tau_r = \tau_0,
\]
where $\tau_r$ is the torsion of $C_r(t)$ and $\tau_0$ is the unconditional prime‑supported distribution defined in (1).
\end{enumerate}
\end{theorem}

\begin{proof}[Sketch of equivalence]
If (i) holds, the choice $r(t)\equiv\frac12$ is smooth and satisfies the interpolation condition.  The curve $C_{1/2}(t)$ is simply a scaling of the unit‑radius helix by the constant factor $\frac12$, and scaling does not change the torsion.  Hence $\tau_{1/2} = \tau_1 = \tau_0$ in the distributional limit, proving (ii).

Conversely, assume (ii).  The torsion $\tau_r$ consists of the contribution from the angular part (which yields $\tau_0$ when $r$ is constant) plus extra singular terms coming from $r',r''$.  A direct calculation shows that those extra terms are linear combinations of $\delta$‑derivatives supported at the same points $t=m\log p$, with coefficients involving $r'$ and $r''$.  For the sum to equal the pure $\delta$‑measure $\tau_0$ (with no derivatives of $\delta$), all coefficients of higher singularities must vanish.  This forces $r'(m\log p)=0$ and $r''(m\log p)=0$ for all prime powers $m\log p$ that appear.  Since the set $\{m\log p\}$ is dense in $[0,\infty)$ (for large $p$), smoothness of $r$ implies $r'\equiv0$ and $r''\equiv0$ everywhere.  Hence $r$ is constant.  The interpolation condition $r(\gamma_n)=\beta_n$ then forces all $\beta_n$ to equal that common value; by the symmetry $\beta_n\leftrightarrow1-\beta_n$ from the functional equation, the only possibility is $\frac12$.  Thus (i) holds.
\end{proof}

The crucial step in the “(ii)$\Rightarrow$(i)” direction is establishing that the only way the torsion of $C_r$ can reduce to the pure Dirac comb $\tau_0$ is for $r$ to be constant.  This reduces the Riemann Hypothesis to a purely differential‑geometric problem about a specific class of curves with prescribed singularities in their torsion.

\end{document}